\pagebreak

\section{Summary and Conclusion}
We introduce a guideline for documenting AI applications along the entire developing process that extends previous work on model and data cards. For this purpose, we held interviews with certification experts, developers of AI applications preparing for the upcoming regulation of the European Union, and, individuals affected by the introduction of an AI system. All their feedback was incorporated into our framework. 
%We conducted interviews with certification \& standardization bodies and individuals being affected by the introduction of an AI system. Their feedback was merged into a guideline that covers the entire AI development process and takes into account recent research as well as regulation efforts by the EU. 
In contrast to works on model and data cards, we cover the entire development process, which results in defining two new cards: use case card and operation card. Further, we add a perspective on regulation. Our work is also different from most regulation and certification efforts by being more specific and following the development process of AI applications instead of risk dimensions.  Our guideline is meant to support the development of trustworthy and easy to audit AI systems from the beginning and should lay the foundations for a future certification of a system. Although AI regulation is still in an early phase and many changes may be expected, the main foundations and associated risks seem to have emerged. Hence, this work can be seen as a snapshot of the current state and will evolve over time but should already position practitioners well when legal requirements come into force.

%\subsection{Next Steps}
While our guideline is more specific than many previous works, both the field of AI but also the regulation are progressing fast. Hence, certain changes or specifications can be expected. For example, best practices for creating datasets are likely to change in future. For this purpose, a structured and regularly updated collection of best practices and safety measures applied by different organizations similar to the \href{https://incidentdatabase.ai/}{AI incidents database} or the \href{https://appliedaiinitiative.notion.site/Risk-Classification-Database-2b58830bb7f54c9d8c869d37bdb27709}{Risk-Classification Database} could be collected. 



